\documentclass[10pt,a4paper]{amsart}
\usepackage[margin=19mm]{geometry}
\usepackage{amsmath,amssymb,mathtools}
\usepackage[scr=rsfs,cal=euler]{mathalfa}
\usepackage{xfrac}
\usepackage[unicode,hidelinks,bookmarks=false]{hyperref}
\newcommand{\ie}{i.e.\ }
\newcommand{\cf}{cf.\ }
\newcommand{\resp}{respectively\ }
\newcommand{\Subsection}{\subsection}
\providecommand{\nobreakdash}{\nobreak\hbox{-}}
\setlength{\emergencystretch}{2em}
\multlinegap0pt
\usepackage{mathtools}
\usepackage{amssymb}
\newtheorem{thm}{}
\newtheorem{lem}{}
\title{The convergence and uniqueness of a discrete-time nonlinear Markov chain}
\author{Ruowei Li, Florentin M\"unch}
\date{Source: arXiv:2407.00314v2 (CC BY 4.0)}
\begin{document}
\maketitle

\noindent\emph{Source and licence.} The convergence and uniqueness of a discrete-time nonlinear Markov chain. Version arXiv:2407.00314v2 (CC BY 4.0), distributed by arXiv under the Creative Commons Attribution 4.0 International licence, http://creativecommons.org/licenses/by/4.0/, as recorded in the arXiv licence metadata for this exact version and shown on the abstract page https://arxiv.org/abs/2407.00314v2. Full licence notice: \texttt{licenses/arxiv-2407.00314-CC-BY-4.0.txt}.\par\smallskip

\noindent\textit{Editorial scope.} Counted unit: the article's thm 3 (source0.tex lines 258--273). The article's complete original proof of it is delivered with it (source0.tex lines 702--806, one \texttt{proof} environment of 105 lines, which the article places away from the statement). No mathematics is written, repaired or re-derived: the statement and the proof are the article's own lines, in the article's own notation, and no retained line is substituted or re-flowed.  Retained uncounted as the source-local context the proof needs: lines 221--234; lines 671--680.  Nothing was omitted from any retained span, so the package carries no omission record.  No retained line contains a citation, so the wrapper carries no bibliography and no bibliography entry.  No retained line contains a figure environment, an image-inclusion command or drawing code, so no image is delivered and the package declares no asset dependency.  Editorial formatting only, in addition: an AMS \texttt{amsart} preamble that restates the article's own declarations and macros used by the retained lines, namely the environments \texttt{thm}, \texttt{lem} on separate counters, together with the standard packages those lines need.  The retained environments print in this document as Theorem 1, Lemma 1 in order of appearance, whereas the article prints the same items with the numbering of its own full document.  The complete native source of the article as distributed in the arXiv e-print for this exact version is bundled unchanged as \texttt{sources/161551/source0.tex}.
\begin{thm}
\label{thm:NLMC convergence 2}Let $f\in \mathbb{R}^{N}$. If a discrete-time nonlinear Markov
chain $P:\mathbb{R}^{N}\to\mathbb{R}^{N}$ satisfies

(1) monotonicity,

(2) strict monotonicity of corresponding components,

(5) non-expansion,

(9) $P^{n}f$ has a finite accumulation point $g\in \mathbb{R}^{N}$,

then $Pg=g$ and $P^{n}f\to g$ as $n\to\infty$. 
\end{thm}

\begin{lem}
\label{lem: convergence lemma by the uniqueness}Let $(X,d)$ be a
locally compact metric space. If a sequence $\left\{ x_{n}\right\}\subset X$
has exactly one accumulation point and satisfies 
\[
d(x_{k},x_{k+1})\leq C \quad \text{for some }C>0 \text{ and all } k \in\mathbb{N}_{+},
\]
then $\left\{ x_{n}\right\} $
converges. 
\end{lem}

\begin{thm}
\label{thm:NLMC convergence 3}Let $f\in \mathbb{R}^{N}$. If a discrete-time nonlinear Markov
chain $P:\mathbb{R}^{N}\to\mathbb{R}^{N}$ satisfies

(1) monotonicity,

(5) non-expansion,

(7) uniform connectedness,

(8) accumulation point at infinity, i.e., $f_{n}\coloneqq P^{n}f-P^{n}f(x_{0})\cdot \overrightarrow{1}$
has a finite accumulation point $g\in \mathbb{R}^{N}$, 

then $f_{n}\to g$ as $n\to\infty$
and the limit is unique. 
\end{thm}

\begin{proof}[Proof of Theorem \ref{thm:NLMC convergence 3}]
Let $\tilde{P}=P^{n_{0}}$, then $\tilde{P}$ is uniformly connected
with $n_{0}=1$, implying strict monotonicity of corresponding components (2). As $g$ is an accumulation
point, then there exist $0\leq i\leq n_{0}-1$ and a subsequence $\left\{ n_{k}=m_{k}n_{0}+i\right\} $
such that $f_{n_{k}}\to g$ and $P^{n_{k}}f(x_{0})\to a\in[-\infty,+\infty]$
as $k\to+\infty$. Divide it into two cases: $\mid a\mid<+\infty$
and $\mid a\mid=+\infty$.

\textbf{Case 1.} If $\mid a\mid<+\infty$, w.l.o.g, suppose $a>0$.

Then as $k\to+\infty$, $$P^{n_{k}}f=\tilde{P}^{m_{k}}P^{i}f\eqqcolon\tilde{P}^{m_{k}}F\to g+a\cdot \overrightarrow{1}\eqqcolon\tilde{g}.$$
By Theorem \ref{thm:NLMC convergence 2}, we have
$\tilde{P}^{m}F\to\tilde{g}$ as $m\to+\infty$ and $\tilde{P}^{k}\tilde{g}=\tilde{g}$
for all $k\in\mathbb{N}_{+}$. Then for different accumulation points of $f_{n}$,
by the non-expansion property (5), their corresponding $a$ must be
the same case.

If there are two accumulation points $\tilde{g}_{1}\neq\tilde{g}_{2}$,
suppose $\alpha\coloneqq\underset{1\leq x\leq N}{\text{max}}\left(\tilde{g}_{1}(x)-\tilde{g}_{2}(x)\right)$ $>0$.
Then $\tilde{g}_{1}\leq\tilde{g}_{2}+\alpha \cdot \overrightarrow{1}$. If there exists $x$
such that $\tilde{g}_{1}(x)<\tilde{g}_{2}(x)+\alpha$, then by the
uniform connectedness and non-expansion property of $\tilde{P}$, 
\[
\tilde{g}_{1}=\tilde{P}\tilde{g}_{1}<\tilde{P}\left(\tilde{g}_{2}+\alpha \cdot \overrightarrow{1} \right)\leq\tilde{P}\tilde{g}_{2}+\alpha \cdot \overrightarrow{1}=\tilde{g}_{2}+\alpha \cdot \overrightarrow{1},
\]
which implies $\tilde{g}_{1}<\tilde{g}_{2}+\alpha \cdot \overrightarrow{1}$ and contradicts
to the definition of $\alpha$. Hence $\tilde{g}_{1}=\tilde{g}_{2}+\alpha \cdot \overrightarrow{1}$,
which means $\tilde{g}_{1}-\tilde{g}_{1}(x_{0}) \cdot \overrightarrow{1}=\tilde{g}_{2}-\tilde{g}_{2}(x_{0}) \cdot \overrightarrow{1}$.
Then the two accumulation points $g_{i}=\tilde{g}_{i}-\tilde{g}_{i}(x_{0}) \cdot \overrightarrow{1}$
for $i=1,2$ of $f_{n}$ are the same. By Lemma \ref{lem: convergence lemma by the uniqueness},
we get the convergence of $f_{n}$.

\textbf{Case 2.} If $|a|=+\infty$, w.l.o.g., suppose $a=+\infty$.

Since $\tilde{P}(f+r\cdot \overrightarrow{1})-r \cdot \overrightarrow{1}$ is decreasing for positive $r$ by the non-expansion
condition (5), we define 
\[
Qf\coloneqq\underset{r\to+\infty}{\text{lim}}\left(\tilde{P}(f+r\cdot \overrightarrow{1})-r \cdot \overrightarrow{1}\right).
\]
As $$Qg-g=\underset{k\to+\infty}{\text{lim}}\left(\tilde{P}^{m_{k}+1}F-\tilde{P}^{m_{k}}F\right)$$
is finite by the non-expansion property (5) of $\tilde{P}$, then for
any $f\in\mathbb{R}^{N}$, we know $Qf$ is also finite by the non-expansion property.
Then $Q$ satisfies the non-expansion property (5) and constant additivity
$Q(f+c \cdot \overrightarrow{1})=Qf+c \cdot \overrightarrow{1} $.

Now we show the connectedness of $Q$. For some component $x$, positive
$\delta$ and $f,h\in\mathbb{R}^{N}$ with $f\geq h+\delta\cdot1_{x}$,
since $\tilde{P}$ is uniformly connected with $n_{0}=1$, we have
\[
\tilde{P}\left(f+r \cdot \overrightarrow{1}\right)-r \cdot \overrightarrow{1}>\tilde{P}\left(h+r \cdot \overrightarrow{1}\right)-r\cdot \overrightarrow{1}+\epsilon_{0}\delta \cdot \overrightarrow{1}\geq Qh+\epsilon_{0}\delta \cdot \overrightarrow{1}.
\]
Let $r\to+\infty$, then $Qf>Qh+\epsilon_{0}\delta \cdot \overrightarrow{1}$ and $Q$ is
uniformly connected with $n_{0}=1$ and $\epsilon_{0}$.

Let $\lambda_{+}^{\tilde{P}}(f)\coloneqq\underset{1\leq x\leq N}{\text{max}}\left(\tilde{P}f(x)-f(x)\right)$
and $F=P^{i}f$, then $\lambda_{+}^{\tilde{P}}(\tilde{P}^{m}F)>0$
for all $m\in\mathbb{N}_{+}$ since $\tilde{P}^{m_{k}}F(x_{0})\to+\infty$. By the
monotonicity and non-expansion property, 
\[
\tilde{P}^{m+2}F\leq\tilde{P}\left(\tilde{P}^{m}F+\lambda_{+}^{\tilde{P}}(\tilde{P}^{m}F) \cdot \overrightarrow{1}\right)\leq\tilde{P}^{m+1}F+\lambda_{+}^{\tilde{P}}(\tilde{P}^{m}F) \cdot \overrightarrow{1},
\]
which means $\lambda_{+}^{\tilde{P}}(\tilde{P}^{m}F)$ is decreasing
in $m$. Since $g$ is a finite accumulation point, for all $l\in\mathbb{N}_{+}$,
we have 
\[
Qg-g=\underset{k\to+\infty}{\text{lim}}\left(\tilde{P}^{m_{k}+1}F-\tilde{P}^{m_{k}}F\right)
\]
and 
\[
Q^{l+1}g-Q^{l}g=\underset{k\to+\infty}{\text{lim}}\left(\tilde{P}^{m_{k}+l+1}F-\tilde{P}^{m_{k}+l}F\right).
\]
Then by the monotonicity of $\lambda_{+}^{\tilde{P}}(\tilde{P}^{m}F)$,
we have 
\[
\lambda_{+}^{Q}(g)=\underset{1\leq x\leq N}{\text{max}}\left(Qg(x)-g(x)\right)=\underset{m\to+\infty}{\text{lim}}\lambda_{+}^{\tilde{P}}(\tilde{P}^{m}F)=\lambda_{+}^{Q}(Q^{l}g).
\]
Since $Qg\leq g+\lambda_{+}^{Q}(g)\cdot \overrightarrow{1}$, if there is $x$ such that $Qg(x)<g(x)+\lambda_{+}^{Q}(g)$,
then 
\[
Q^{2}g<Qg+\lambda_{+}^{Q}(g)\cdot \overrightarrow{1}=Qg+\lambda_{+}^{Q}(Qg)\cdot \overrightarrow{1}
\]
by the uniform connectedness and constant additivity of $Q$, which
contradicts to the definition of $\lambda_{+}^{Q}(g)$. Hence $Qg=g+\lambda_{+}^{Q}(g)\cdot \overrightarrow{1}$
and $Q^{l}g=g+l\lambda_{+}^{Q}(g)\cdot \overrightarrow{1}$.

Let $f^{(1)},f^{(2)}\in\mathbb{R}^{N}$ and assume $g_{i}$ is an
accumulation point of $f_{n}^{(i)}=P^{n}f^{(i)}-P^{n}f^{(i)}(x_{0})\cdot \overrightarrow{1}$
for $i=1,2$. Then by the non-expansion property 
\[
\left\Vert g_{1}+l\lambda_{+}^{Q}(g_{1})\cdot \overrightarrow{1}-g_{2}-l\lambda_{+}^{Q}(g_{2})\cdot \overrightarrow{1}\right\Vert _{\infty}=\left\Vert Q^{l}g_{1}-Q^{l}g_{2}\right\Vert _{\infty}\leq\left\Vert g_{1}-g_{2}\right\Vert _{\infty},
\]
we know $\lambda_{+}^{Q}(g_{1})=\lambda_{+}^{Q}(g_{2})$. If $g_{1}\neq g_{2}$,
suppose $\alpha\coloneqq\underset{x}{\text{max}}\left(g_{1}(x)-g_{2}(x)\right)>0$.
Then $g_{1}\leq g_{2}+\alpha \cdot \overrightarrow{1}$ and $0=g_{1}(x_{0})<g_{2}(x_{0})+\alpha=\alpha$.
By the uniform connectedness of $Q$, 
\[
g_{1}+\lambda_{+}^{Q}(g_{1})\cdot \overrightarrow{1}=Qg_{1}<Qg_{2}+\alpha \cdot \overrightarrow{1}=g_{2}+\lambda_{+}^{Q}(g_{2}) \cdot \overrightarrow{1}+\alpha \cdot \overrightarrow{1},
\]
which implies $g_{1}<g_{2}+\alpha \cdot \overrightarrow{1}$ and contradicts to the definition
of $\alpha$. Then the two accumulation points are the same.

By Lemma \ref{lem: convergence lemma by the uniqueness} and letting
$f^{(1)}=f^{(2)}=f$, we get the convergence of $f_{n}$. Moreover,
the case $f^{(1)}\neq f^{(2)}$ shows the uniqueness of the limit. 
\end{proof}

\end{document}
